\chapter{CON LE MANI IN MANO}\label{con-le-mani-in-mano}

\hfill\break

Intorno a mezzanotte, riuscii ad andare in branda, un letto vero,
abbastanza a lungo da potermi stendere. Dormire in un letto vero era
delizioso, quanto mangiare cibo vero. Tutto quello che volevo era
togliermi gli stivali e strisciare sotto le coperte, ma Skippy voleva
parlare. Di nuovo. «Oh, cavolo, Skippy, puoi lasciarmi dormire un po'?»

«Joe, sono solo. Ecco, l'ho detto.»

«Parlo con te di continuo.»

«Joe, mi sento solo mentre sto parlando con te. Parli così lentamente e
ti ci vuole così tanto tempo per arrivare al punto, che è come se stessi
ad aspettare ogni giorno accanto alla cassetta della posta una lettera
che contiene una sola parola. E dovessi aspettare un altro giorno intero
per la parola successiva. Quando aspetto e finalmente ricevo una lettera
che quel giorno dice solo ``ehm'', ho voglia di urlare. Cazzo! Mi sembra
di dovermi sporgere dentro la tua gola per trascinarti fuori le parole.
Parli così lentamente. Dillo! Dillo! Cazzo! Tira fuori le parole!»

Per un istante, potei sbirciare dentro il vero Skippy, negli eoni di
dolore che aveva sopportato. Se avessero lasciato me tutto solo per un
mese soltanto, il mio sistema diagnostico avrebbe rilevato che ero del
tutto impazzito. Non potevo immaginare come si sentisse: «Hai bisogno di
più persone con cui parlare».

«Joe, un altro paio di persone non farebbero\ldots»

«Skippy, ci sono miliardi di esseri umani su questo pianeta, e tutti
pensano di avere qualcosa di importante da dire. Purtroppo, intasano
internet con blog, vlog, chat, video di gattini e litigano sullo sport.
Senti, non dovrei suggerirtelo, perché la tua esistenza dovrebbe essere
segreta, ma se non dici a nessuno che sei un'intelligenza artificiale
aliena, nessuno si farà male, giusto?»

«Miliardi.»

«Sì, miliardi. Per lo più stupidi come me, ma potresti trovare dei vermi
piatti là fuori, tra i batteri.»

«Dubito. Ci penserò su, grazie.»

Non riuscivo a smettere di sbadigliare. «Io stacco per un po', non
metterti nei guai mentre dormo, ok?»

\hfill\break

Sei ore buone di sonno ininterrotto fecero miracoli per il mio umore. Mi
svegliai quando un sergente in uniforme dell'Aeronautica militare bussò
alla mia porta e mi portò un vassoio con una caraffa di caffè caldo e
una tazza. «Colazione fra trenta minuti. La doccia è in fondo al
corridoio, alla sua sinistra.» E aggiunse: «Signore», anche se poteva
vedere i galloni di rango sulla mia nuova serie di uniformi appese
nell'armadio.

«Oh, questo è un buon caffè.» Il primo sorso sulla lingua fu inebriante.

«Hai dormito bene?», chiese Skippy.

«Sì, grazie, ho russato?»

«No, non molto dai. Sono contento che tu abbia dormito bene, ci aspetta
una giornata impegnativa.»

Sembrava essere di buon umore, per una solitaria intelligenza
artificiale vecchia milioni di anni. «Come stai? Neanch'io ti ho sentito
russare.»

«Sto bene, grazie», disse allegro. «Bella giornata, eh?»

«Oh, merda.» Il fatto che fosse gentile con me poteva significare una
sola cosa: si era cacciato nei guai. «Che cos'hai combinato ieri sera?»,
chiesi molto lentamente.

«Ho seguito il tuo consiglio e sono andato su internet per incontrare un
po' di persone.»

«Un po'?»

«Finora, un miliardo e centootto milioni, più o meno. Al momento sto
chattando o messaggiando o inviando e-mail con circa trentanove
milioni.»

«Tutti allo stesso tempo?»

«Mi tiene occupato, senza sforzare le mie capacità.»

«Occupato, va bene.»

«Una mente oziosa è il parco giochi del diavolo. L'ho sentito da un
predicatore del fuoco eterno dell'Idaho ieri sera. Cavolo, non
crederesti ai porno che ha sul suo\ldots»

«Non ne ho bisogno! Cazzo, mi sono appena svegliato, Skippy.»

«È probabile che sia meglio così. Alla tua specie piace molto la
pornografia, il che è impressionante, per una specie con due soli
generi.»

«Wow! Siamo i numeri 1!» Sollevai un immaginario dito di schiuma
ingoiando il caffè. «Sì! Gli umani sbaragliano tutti!»

«Le differenze fisiche tra maschi e femmine sono così lievi, che non è
nemmeno\ldots»

«Ah, ma quelle differenze significano tutto. Fidati di me.»

«Ti credo sulla parola», immaginai che stesse strabuzzando gli occhi là
dentro.

«Quindi, hai incontrato delle persone online. E?»

«La maggior parte di loro pensa che sia uno stronzo o l'equivalente
nella loro lingua madre.»

«Sono scioccato!»

«Non fare lo stronzo.»

«Scusa. Ti piace conoscere gente?»

«Mettiamola così, hai mai scorso fino in fondo una pagina web per
leggere i commenti degli utenti?»

«Oh, cavolo, Skippy, mai leggere i commenti, lo sanno tutti! Quei
commenti sono tutti scritti da tizi che se ne stanno seduti in mutande,
perché non hanno nient'altro da fare.»

«Sono d'accordo con te, tranne per la parte sulla biancheria intima. Ho
acceso alcune webcam senza che loro lo sapessero.»

«Bleah.»

«Bleah davvero. Mi sbagliavo, non tutte le scimmie sono glabre. A un
ragazzo sembrava crescesse un tappeto peloso sulla schiena. Questi sono
un paio di petabyte di memoria che vorrei poter cancellare. Ehi, a
proposito di petabyte e di altri grandi numeri\ldots»

«Che mi dici?»

«Beh, eh eh, questa è una storia divertente\ldots»

Ah, oh. «Divertente del tipo ``ah ah'', o divertente del tipo io che
passo del tempo nella prigione federale?»

«Non hai fatto niente.»

«A te non metteranno in prigione. Quindi, cosa hai fatto?»

«Beh, la larghezza di banda internet di qui è striminzita, anche se
comprimo i miei messaggi, così ho raggiunto un posto che sembra avere
connessioni con tutto.»

«Google?», tirai un sospiro di sollievo.

«No, c'è un posto chiamato Fort Meade, nel Maryland? È la vostra Agenzia
per la sicurezza nazionale. Un tizio lì si sta strappando i capelli da
mezzanotte, cercando di capire chi ha hackerato il loro sistema.»

«Skippy!» Mi si gelò il sangue. Mi veniva da vomitare. «Oh, sono in guai
grossi.» Mi stavo già pentendo di avere bevuto il caffè.

«No, siamo a posto, pensano che sia un quindicenne di Fresno di nome
Billy. Gli ho dato una pagina Facebook e una famiglia finta e ho
retrodatato un sacco di file. A quanto pare, c'è una squadra dell'Fbi
diretta da Starbucks, da dove pensano che mi colleghi. Saranno così
delusi. Ehi, sto guardando dalla telecamera di sicurezza del locale,
vuoi vedere?»

«No! Skippy, non puoi farlo.»

«Certo che posso.»

«Volevo dire che non dovresti farlo.»

«Vedi? Le lingue umane sono così imprecise.»

«Non è divertente.» Aprii la porta, certo che guardie armate fino ai
denti stessero venendo ad arrestarmi in quel momento. «L'Agenzia per la
sicurezza nazionale ha dati top secret!»

«Non per me. E poi non m'importa di nessuna delle stronzate che
ritengono tanto importanti.»

«Hai detto qualcosa a qualcuno delle informazioni che hai trovato?»

«No, andiamo, Joe, nessuna delle stronzate segrete del governo vale il
mio tempo. Ma, ehi, vuoi sapere come ha fatto la Nasa a fingere
l'atterraggio sulla luna?»

«Che cosa?»

«È uno scherzo. Sto scherzando.»

«Non scherzare su queste cose. Ehm, stavi scherzando davvero, giusto?»

«Posso mostrarti i primi piani del sito di atterraggio dell'Apollo 11
rilevati dai sensori dell'Olandese Volante, se vuoi. L'ho esaminato
ieri.»

La mia curiosità superò la mia paura. «Che aspetto ha?»

«La tua specie è molto più coraggiosa che intelligente. Per me il
viaggio nello spazio è una passeggiata, ma i tuoi astronauti sono andati
lassù in barattoli di latta. Anche io sono impressionato. Si può quasi
leggere l'etichetta della zuppa Campbell sui loro velivoli
d'atterraggio, accanto al logo della Nasa. Ragazzi, voi scimmie avevate
davvero le palle per atterrare sulla luna con quella tecnologia
schifosa.»

«Ottimo, grazie. Puoi, per favore, per favore, lasciare in pace
l'Agenzia per la sicurezza nazionale? Puoi farlo per me?»

«Perché? Mi sto divertendo. Non mi divertivo da milioni di anni!»

«Ci sono altri modi per divertirsi, Skippy, che non implicano che io
finisca in una prigione federale.»

«Sei già uscito di prigione. Due volte.»

«Puoi essere serio per un minuto? Mi sono arruolato nell'esercito
americano, non posso contribuire a incasinare le infrastrutture di
sicurezza americane.»

«Ma causare problemi è molto divertente», brontolò Skippy.

«Vuoi divertirti a creare problemi? Vai online e fai girare una voce
credibile sul fatto che Justin Bieber interpreterà Darth Vader nel
prossimo film di Star Wars», suggerii.

«Oooh, bella questa! Sapevo che c'era un buon motivo per frequentarti,
Joe. Ok, ho appena messo online quattordici minuti di quello che sembra
un filmato di studio piratato.»

Mi diedi uno schiaffo sulla fronte: «Oh, mio Dio, che cosa ho fatto?».

«E ho pubblicato uno studio del ministero della Salute secondo il quale
Vegemite funziona meglio del Viagra.»

«Smettila! Vado a farmi una doccia, cerca di non tirare bombe atomiche
mentre sono via.»

«Sai che non farei del male a nessuno, colonnello Joe. Mmh, che ne dici
di\ldots»

Chiusi la porta e mi diressi in fondo al corridoio, verso i bagni,
chiedendomi se avessi tempo per comprare azioni dell'azienda produttrice
della crema salata Vegemite, qualunque fosse.

\hfill\break

Non ero mai stato così nervoso in vita mia. Incontrare la presidentessa
all'esterno, dopo essere appena scesi dall'orbita in una navetta
thuranin, con una lattina di birra cromata al seguito, aveva reso
l'intera esperienza così irreale, che avevo dimenticato di essere
nervoso. Adesso che indossavo un'uniforme appena stirata, sedermi in una
stanza con i capi di stato maggiore, i capi della Cia e dell'Agenzia per
la sicurezza nazionale, il direttore della sicurezza nazionale e il
consulente scientifico e il capo di gabinetto della presidentessa mi
fece seccare la bocca e quasi pisciarmi addosso. La stanza era un po'
come lo Studio Ovale, con una grande scrivania a un'estremità, due file
di divani e un tavolino. Mi fecero sedere in fondo a un divano, il più
vicino alla sedia su cui si sarebbe seduta la presidentessa. Il capo di
stato maggiore dell'esercito, che avevo visto solo in foto prima di
allora, era seduto proprio accanto a me. Mentre aspettavamo la
presidentessa, la gente beveva caffè da tazzine di porcellana appoggiate
su piattini, immagino per non lasciare macchie sul tavolo. Le tazzine
non erano quelle con il sigillo presidenziale, dovevano averle lasciate
a Washington. Quando mi fui seduto ed ebbi il tempo di guardarmi
attorno, la stanza non sembrava affatto maestosa come lo Studio Ovale.
Infatti, pensai che i mobili potevano benissimo essere stati presi dalla
hall del Ramada medio. Oppure, dato che poi mi resi conto che il braccio
del divano era consumato e aveva una chiazza misteriosa, forse da un
Motel 6. Il capo dell'esercito prese la sua tazzina di caffè in una mano
abbastanza grande da coprirla per intero e io cercai di imitarlo. La
tazzina tremò così tanto che fece un gran fracasso urtando il piattino,
perciò l'appoggiai con l'attenzione che si impiegherebbe per
disinnescare una testata nucleare. Il capo dell'esercito, il generale
Brenner, ebbe pietà di me: «Meglio se lasci stare il caffè», disse
calmo.

Dovetti deglutire due volte per avere abbastanza umidità in gola da
rispondere. «Signore, sono stato in combattimento e non ero così
nervoso.»

Il generale, che era stato in un sacco di combattimenti, annuì: «Alzati
quando arriva la presidentessa, e non rovesciare il tavolo. Parla quando
ti parlano e, quando parli, sii diretto. Non abbiamo tempo per le
stronzate».

\hfill\break

La porta si aprì, un paio di agenti dei servizi segreti entrarono,
seguiti dalla presidentessa. Sembrava fosse stata in piedi quasi tutta
la notte.

Il direttore della sicurezza nazionale fu il primo a parlare: «Signora
presidentessa, ho un aggiornamento sulla violazione della sicurezza
all'Agenzia per la sicurezza nazionale di ieri sera».

Sbiancai come un cencio. Speravo che nessuno mi stesse guardando.

Il tizio continuò: «Abbiamo contenuto la situazione e stiamo valutando
il danno ora. La persona che pensavamo fosse il colpevole si è rivelata
essere un fantasma, una falsa pista. Questo è stato un attacco molto
sofisticato, è probabile che i kristang\ldots».

«Kristang? Ah ah, quegli idioti?! No! Ero io», disse una voce ovattata
dal mio zPhone.

Ora tutti stavano fissando me. «Non io, io.» Tirai fuori il dispositivo
dalla tasca e lo appoggiai sul tavolo.

«Sì. Ero io, Skippy il Magnifico. Ehi, come mai non sono stato invitato
a questa festa? Sembra una figata.»

«Se non ho capito male, l'essere di nome Skippy», il direttore della
sicurezza nazionale mi rivolse uno sguardo ostile, «è entrato
nell'Agenzia per la sicurezza nazionale ieri sera e ha saccheggiato i
nostri file top secret?»

«Saccheggiato? Ho lasciato tutto dov'era. Se non volevate che la gente
leggesse i vostri file, dovevate criptare i dati», brontolò Skippy.

«Sono criptati!»

«Davvero? Oh, pensavo che quei file fossero solo mal indicizzati.
Ehm\ldots{} Era crittografia? Oh, ah ah. Mi stai prendendo in giro,
giusto? Buona questa.»

«Oh, Cristo!» Il consigliere scientifico nazionale sussultò: «Quello è
lo stato dell'arte della crittografia! Come hai ottenuto le chiavi?».

«Chiavi?», chiese Skippy con fare innocente.

Decisi di porre fine al divertimento della mia lattina di birra:
«Signore, col nostro basso livello di tecnologia, Skippy non si prende
la briga di decriptare i file, gli basta saltare alla fine per leggere
il contenuto. Ha qualcosa a che fare con il gatto di Schröder\ldots»,
alzai le mani.

«Il gatto di Schrödinger. Comunque, che si dice, amici?», incalzò
Skippy.

La presidentessa sorrise: «Signor Skippy, visto che non possiamo
escluderla dalla riunione, vuole unirsi a noi?».

«No, va bene così, stavolta starò al telefono. In questo modo posso
starmene sdraiato qui sul divano, con le pantofole pelose e le mutande e
fingere di ascoltare. Inoltre, sto guardando la Ruota della Fortuna.» E
probabilmente tutti gli altri programmi televisivi trasmessi in tutto il
pianeta in quel momento.

Alzai di nuovo le mani.

«Ho visto», annunciò Skippy.

«Come?», si accigliò il consulente scientifico.

«Attraverso la fotocamera del telefono di Joe. Ma tu pensa. Inoltre, le
particelle di polvere nell'aria contengono ioni che, mmh, meglio che non
ne parli con voi scimmie. È tutto troppo complicato.»

«Scimmie? Cosa\ldots», iniziò a dire il capo di stato maggiore
dell'Aeronautica militare.

«Molto bene», disse la presidentessa. «Signor Skippy, abbiamo monitorato
i siti kristang sulla Terra, può dirci lo stato della loro nave in
orbita?»

«Certo. Ci sono due kristang ancora vivi a bordo della nave da trasporto
truppe. Altri due sono sopravvissuti all'inizio, perché erano in
compartimenti che potevano sigillare a mano. Ma il loro ossigeno si è
esaurito, quindi sono morti. Gli altri due erano già in tuta spaziale e
si stavano preparando per uscire, quando i loro amici sono stati
soffiati nello spazio. Da me. Quei due sono riusciti a sigillare le
porte, ripristinare l'atmosfera di una parte della nave, e stanno
tentando di ottenere l'accesso al compartimento di stoccaggio delle armi
biologiche, in modo da poter lanciare missili contro di voi.»

«Armi biologiche?», esclamò allarmato il capo di stato maggiore della
Marina. «Che tipo di armi biologiche?»

«Oh, niente di speciale. Virus modificati aerosolizzati, armi
geneticamente modificate basate sui comuni raffreddore, influenza, virus
Ebola e febbre di Marburg. Non sono ancora molto efficienti, perché i
kristang non hanno avuto granché tempo per studiare la vostra biologia.
I test che hanno fatto su Campo Alfa indicano una letalità nella prima
settimana di appena il 12\%, ma la letalità sale al 62\% entro un mese.
È il fatto che il sistema immunitario del soggetto venga attaccato da
più virus allo stesso tempo che logora le persone e le uccide», disse
Skippy in tono molto oggettivo.

La stanza era tutta in tumulto, a parte me che, seguendo il consiglio
del capo dell'esercito, tenevo la bocca chiusa. La presidentessa alzò di
nuovo le mani per riportare il silenzio: «Signor Skippy, la prego, ci
dica\ldots».

«So cosa mi stai per chiedere ora, quindi eccolo qui. I test su Campo
Alfa sono stati effettuati su soggetti umani, catturati sulla Terra e
introdotti lì di nascosto, non sul personale militare assegnato
d'ufficio. Il personale militare è in prevalenza maschile, inoltre tende
a essere più giovane e più in forma rispetto alla media della
popolazione umana in generale, quindi non sarebbe un buon campione
rappresentativo per i test sulle armi biologiche. I kristang hanno
rapito una sezione incrociata di età, generi e gruppi etnici, e l'hanno
portata dalla parte opposta rispetto alle strutture militari del pianeta
Alfa.»

«È terrificante», commentò la presidentessa in tono calmo, e nessun
altro parlò.

«Le armi biologiche a bordo della nave da trasporto truppe in orbita
contengono abbastanza virus aviotrasportati da uccidere diversi milioni
di esseri umani nell'onda iniziale. C'è solo una fornitura limitata di
armi biologiche, ma quei missili sono mirati ai principali centri
abitati come San Paolo, Shanghai, Tokyo, Mumbai, New York, tutti i
soliti sospetti. L'uso di armi biologiche è decisamente contro Le
Regole, ma i kristang qui pensano che la Terra sia così lontana dalla
civiltà che ne varrebbe la pena, se la loro situazione si facesse
disperata. Come lo è ora.»

«In che modo possiamo impedire ai kristang lassù di lanciare quelle
armi? Abbiamo ancora missili nucleari», dichiarò il direttore della
sicurezza nazionale. «Quella nave è in un'orbita troppo alta per\ldots»

«Oh, non ce n'è bisogno», intervenne Skippy in tono allegro, «ho
sterilizzato tutto il contenuto delle armi biologiche a bordo della nave
e disattivato quei missili ieri. Ah, ehi, è probabile che avrei dovuto
dirvelo prima, eh?» Mi diedi uno schiaffo sulla fronte mentre parlava.
«Cercare di lanciare quelle armi sta tenendo occupati i due kristang, e
finché staranno felicemente cercando di sterminare la vostra
popolazione, non creeranno alcun problema, quindi li lascerò fare. A un
certo punto, dovrete mandare un commando o qualcosa del genere lassù per
ucciderli, perché quei due potrebbero annoiarsi e provare a
sovraccaricare il loro reattore a fusione. Il che mi ricorda che io,
ehm, ok, uhm, ho appena iniziato a spegnere il loro reattore. Mi occupo
anche di quel problema.»

Le facce nella stanza sbiancarono, mentre il sangue rifluiva sotto le
suole. Mi nascosi il volto tra le mani. «Skippy», chiesi, «come fa un
essere d'incredibile intelligenza a essere tanto distratto?»

«Vedi?», replicò Skippy con fare innocente. «Questo è il genere di cose
che dovresti ricordarmi tu, colonnello Joe. Non riesco a pensare a
tutto.»

«Tute spaziali», borbottai a mezza voce.

«Oh, stai zitto, scimmione.»

«Tute spaziali?», chiese la presidentessa.

«È una lunga storia», disse Skippy.

La presidentessa lanciò un'occhiata al direttore della sicurezza
nazionale. «Mmh mmh. Tutto sembra essere una lunga storia con voi due.»

«Skippy», mi sentivo ancora un idiota a chiamarlo così, di fronte ai
vertici della nazione, «possiamo manipolare da remoto i robot a bordo
dell'Olandese, per occuparci di quei due kristang?» Non potevo
immaginare un commando umano, con indosso tute spaziali della Nasa,
impegnato in un confronto a fuoco con un paio di guerrieri kristang, a
bordo di una nave kristang.

«Certo. Sei pieno di buone idee, colonnello Joe.»

Il capo di stato maggiore dell'esercito si girò verso di me: «Dovremo
portare alcuni ranger sull'Olandese, per usare questo equipaggiamento di
manipolazione da remoto. Serve a controllare a distanza i robot da
combattimento thuranin, giusto?».

«Possiamo fornire una squadra di forze speciali Seal», offrì il capo
della Marina, per non essere da meno.

«Ehm, sì, signori, ma sarebbe meglio mandare alcuni dei nostri pirati,
voglio dire, alcuni della ciurma dell'Olandese lassù. Abbiamo esperienza
nel controllo di quei robot in combattimento. Le abilità da videogioco
sono più utili del tipo di cose per cui i ranger o gli uomini del Seal
sono addestrati, signori.» Ero sicuro che, in qualche modo, avremmo
avuto ranger dell'esercito, uomini delle forze speciali Seal e del Force
Recon del corpo dei Marine e ufficiali di tattica speciale
dell'Aeronautica a bordo dell'Olandese quando sarebbe andata in orbita.
Con tutta probabilità anche dell'Hostage Rescue Team del Fbi, dato che
nessuno voleva essere escluso dall'azione. E quelli erano solo gli
americani. Si prevedeva un grande affollamento. Avremmo dovuto portare
con noi un sacco di deodoranti ambientali.

Avrei chiesto a Skippy di prenderne nota.

\hfill\break

La riunione si trascinò, ciascuno dei consiglieri della presidentessa
fece rapporto sulla rispettiva area di responsabilità, fino a quando non
arrivammo alla questione che stava a cuore a tutti: i kristang sulla
Terra. Erano sopravvissuti perché si trovavano in bunker sotterranei che
non potevano essere danneggiati dalle armi che Skippy aveva selezionato
nel nostro attacco iniziale. Le uniche armi in grado di raggiungerli
erano bombe atomiche di costruzione umana o il cannone a rotaia a bordo
dell'Olandese. Una sola testata nucleare tattica non avrebbe avuto la
potenza sufficiente per arrivare abbastanza lontano; secondo le stime di
Skippy, avremmo avuto bisogno di scavare un buco profondo, far cadere la
testata nucleare e farla detonare sopra il bunker delle lucertole.
Avevamo moltissimi mezzi per il traforo, anche se ci sarebbero voluti
mesi per posizionare l'attrezzatura e scavare tunnel fino a quel punto.
Skippy avvertì che l'uso di testate nucleari su un mondo abitabile,
inclusa la Terra, era comunque contro Le Regole, regole cui l'umanità
era ora vincolata, poiché avevamo partecipato alla guerra. Se una delle
due parti in guerra avesse mai raggiunto la Terra e scoperto che gli
umani avevano usato delle testate nucleari contro i kristang, le
conseguenze sarebbero state disastrose, come se non avessimo già
abbastanza problemi. Il cannone a rotaia dell'Olandese era l'unica
opzione praticabile. Il che sarebbe stato un problema in due dei siti
interessati, uno vicino alla città di Lione, in Francia, l'altro appena
a ovest di Hangzhou, in Cina. Non avremmo assolutamente potuto usare le
testate nucleari così vicino a quelle città in ogni caso. Il terzo
luogo, la loro base principale, si trovava sotto una montagna a
nord-ovest di Durango, in Colorado.

Era chiaro che alla presidentessa non piaceva l'idea di raccomandare ai
governi francese e cinese di evacuare le loro città in modo che Skippy
potesse usare i penetratori del cannone a rotaia come insetticida.
«Signor Skippy, c'è qualche possibilità che i kristang, dopo un po', si
arrendano? Devono sapere che la loro situazione è senza speranza.»

«No, non lo sanno», spiegò con pazienza Skippy. «Ho interrotto le loro
comunicazioni e cortocircuitato la maggior parte delle loro
apparecchiature elettroniche. Tutto quello che sanno è che una nave
madre thuranin è saltata in orbita, la loro fregata è saltata via e sono
stati attaccati dai loro stessi missili e satelliti. Dal loro punto di
vista, questa è una disputa commerciale tra il clan Vento Bianco dei
kristang e i thuranin, che potrebbe essere stata provocata dal semplice
fatto che i patroni si sono arrabbiati per i ritardi nei pagamenti dei
servizi di spedizione. È già successo prima. Non a questo livello, ma i
kristang non se ne stupirebbero molto. È probabile che le lucertole che
si nascondono nei loro buchi contino di rimanerci per un po', finché,
pensano, i thuranin non avranno capito di avere raggiunto il loro scopo,
o si annoieranno e se ne andranno. Poi le lucertole potranno uscire e
riprendere a essere pessimi ospiti.»

«Cazzo, non funzionerà», ringhiò il generale Brenner. «In qualche modo
possiamo parlare con loro, mostrare loro in che casino si trovano?»

«Certo, se vuole. Parlare con loro non cambierà niente, nessun kristang
si arrenderebbe a una specie primitiva come quella umana, sarebbe
un'umiliazione inimmaginabile. L'unico motivo per cui si arrenderebbero
sarebbe la speranza che le navi kristang possano tornare qui un giorno,
ma, se ciò accadesse, ogni kristang che si fosse arreso agli umani
verrebbe giustiziato, e le loro famiglie a casa verrebbero punite con
severità.»

«Va bene», disse Brenner, «questi proiettili di cannone a rotaia, che
cosa comporterebbe il loro uso?»

«Il cannone a rotaia sull'Olandese Volantenon è stato progettato per il
bombardamento orbitale, perciò, per ottenere la potenza sufficiente a
raggiungere i kristang in profondità nei loro nascondigli, bisognerebbe
caricarlo fino alla potenza massima, cioè circa quaranta minuti tra un
colpo e l'altro. I siti di Lione e Hangzhou non sono molto profondi, due
colpi su ogni sito dovrebbero penetrare abbastanza in profondità da
uccidere tutti i kristang in quei due bunker. Il sito di Durango
richiederebbe tre colpi. A Durango, i proiettili farebbero crollare la
montagna sopra il bunker, chiudendoci dentro le lucertole per sempre.
Per scendere fino al bunker, ci vorrebbero una mezza dozzina di colpi,
che rimuoverebbero la maggior parte della montagna. Cosa che non credo
vogliate, e non è necessaria. Inoltre, troppi attacchi a ipervelocità
getterebbero una notevole quantità di polvere nella vostra atmosfera e
provocherebbero un'alterazione temporanea del clima, come a seguito
dell'eruzione di un vulcano. Tre colpi dovrebbero risolvere il problema
a Durango. Ogni impattatore produce settanta chilotoni in una carica
cava, l'effetto dell'esplosione sarebbe per lo più focalizzato verso il
basso, ma una grande quantità di detriti verrebbe gettata fuori in uno
schema a spruzzo. Lione sarebbe colpita più di Hangzhou.»

La presidentessa contrasse le labbra: «Devo pensarci». Si rivolse al
direttore del Fema, l'ente generale per la gestione delle emergenze: «In
ogni caso, iniziare a evacuare chiunque nel raggio di ottanta
chilometr\ldots».

Skippy emise un convincente colpo di tosse: «Centosessanta sarebbe
meglio. E ci sarà un sacco di polvere e detriti in aria sottovento».

«Nel raggio di centosessanta chilometri da Durango», continuò la
presidentessa.

«Sì, signora», annuì il direttore del Fema, «la maggior parte della
gente si è allontanata da quell'area dopo che le lucertole ci si sono
trasferite.» Il tipo sembrava del tutto esausto, dubitai che avesse
chiuso occhio la notte precedente. O durante molte delle notti e dei
mesi precedenti. Persino la presidentessa aveva dei cerchi scuri sotto
gli occhi e sembrava molto più vecchia di come la ricordavo prima del
Giorno di Colombo. Per quanto le nostre condizioni si fossero
deteriorate su Paradiso, la vita sulla Terra era stata ben peggiore.

Ora mi sentivo in colpa per aver trascorso sei ore di beato riposo nel
mondo dei sogni, mentre tutti gli altri avevano lavorato.